\section*{3 Preliminaries (author source p.11:5)}
We assume that the reader is familiar with the notion of a Boolean circuit and the associated function that it computes. We emphasize that, unless otherwise stated, throughout this work by ``circuit'' we refer to constant fan-in circuits over the standard De Morgan base.
\subsection*{3.1 Uniformity}
We say that a family $\mathcal C=\{C_n\}_{n\in\ZZ^+}$ of Boolean circuits is log-space uniform if there exists a Turing machine that on input $1^n$ uses $O(\log n)$ space and prints the description of the circuit $C_n$, on a write-only output tape. The description is a listing of the gates in an (arbitrary) topological ordering along with a specification of the gate operation and pointers to each of its inputs. In Appendix B we also discuss an extension of our main result for linear-time uniformity.
\subsection*{3.2 Boolean Circuits for Uniselection}
The following standard proposition gives a linear-size circuit, that given a string $x\in\Sigma^n$, over an alphabet $\Sigma$, and an index $i\in[n]$, outputs the symbol $x_i$ (in other words, a multiplexer).
\begin{originalresult}{Proposition 2}
Let $s=s(n)\in\ZZ^+$ and let $\Sigma=\bits^s$. There exists a log-space uniform Boolean circuit for $\Sel_\Sigma^{n\to1}$ with size $O(n\cdot s)$ and depth $O(\log n)$.
\end{originalresult}
\begin{proof}
Without loss of generality assume that $s=1$ (for $s>1$ we can use $s$ parallel copies of the circuit for $s=1$). The circuit follows a divide and conquer strategy. Assume for simplicity that $n$ is a power of 2. Let $x\in\bits^n$ and $i\in[n]$ be the inputs. Let $i_1$ denote the most significant bit of $i$ and let $i'$ denote the remaining bits (i.e., $i'\in[n/2]$). To select the $i$th bit of $x$, the circuit recursively selects the $(i')$th bit of both the lower and upper halves of $x$. Then, based on $i_1$, it decides which of the two bits to output. Overall the circuit size satisfies the recursion: $S(n)=2S(n/2)+O(1)$ and the depth satisfies $D(n)=D(n/2)+O(1)$. The proposition follows.
\end{proof}
\subsection*{3.3 Small-Size Sorting Circuits}
Our circuits for multiselection heavily rely on circuits for sorting integers.
\begin{originalresult}{Lemma 3}
For every $k=k(n)$ and $m=m(n)$, there is a log-space uniform Boolean circuit of size $O(n\cdot m\cdot\log n)$ for sorting $n$ integers of $m$ bits.
\end{originalresult}
This lemma follows immediately from the sorting networks of Ajtai, Koml\'os, and Szemer\'edi [1]. In order to make our multiselection circuit have logarithmic depth, we need a refined version of Lemma 3, which we present in Section 4.
